\chapter*{Sammendrag}

\begin{table}[ht]
  \centering
  \renewcommand{\arraystretch}{1.3}
  \begin{tabularx}{\textwidth}{|X|X|X|}
    \hline
    \multicolumn{2}{|l|}{\textbf{Title:}} & \textbf{Date:} \\ 
    \multicolumn{2}{|l|}{\makecell[l]{Tømmertransport -- utvikling av nye verktøy \\ for varierende vær og føre}} & 20.05.2025 \\ \hline

    \multicolumn{3}{|l|}{\textbf{Deltakere:}} \\ 
    \multicolumn{3}{|l|}{Erik Bjørnsen, Simon Husås Houmb} \\ \hline

    \multicolumn{3}{|l|}{\textbf{Veileder(e):}} \\ 
    \multicolumn{3}{|l|}{Peter Stefan Nussbaum} \\ \hline

    \multicolumn{3}{|l|}{\textbf{Oppdragsgiver:}} \\ 
    \multicolumn{3}{|l|}{Skogkurs} \\ \hline

    \multicolumn{3}{|l|}{\textbf{Nøkkelord:}} \\ 
    \multicolumn{3}{|l|}{Kartløsning, GIS, Fullstack, Skogbruk, Bæreevne} \\ \hline

    \textbf{Antall sider:} & \textbf{Antall vedlegg:} & \textbf{Publiseringsavtale:} \\ 
    \textbf{\pageref{LastMainPage}} & \textbf{\total{chapter}} & \textbf{Åpen} \\ \hline

    \multicolumn{3}{|l|}{\textbf{Kort beskrivelse av bacheloroppgaven:}} \\
    \multicolumn{3}{|p{0.96\textwidth}|}{
    Klimaendringene har forlenget den snøfrie sesongen og økt variasjonen i skogsbilveienes tilstand på grunn av vekslende perioder med frost, tørke og nedbør. Skogkurs har sett et behov for et digitalt verktøy for å vurdere og forutsi bæreevnen til veiene. Denne oppgaven presenterer et \textbf{prototypesystem for heldigital modellering av skogsbilveiers bæreevne under varierende forhold gjennom året}. Systemet er implementert som en webapplikasjon som visualiserer bæreevnen til norske skogsbilveier basert på miljøfaktorer som jordfuktighet, teledybde og løsmasser. Løsningen klassifiserer veiforholdene ved hjelp av en \textbf{trafikklys-modell} og integrerer geo- og meteorologiske data for å vise både nåværende og forventet bæreevne opptil ni dager frem i tid. Systemet består av en backend med en RESTful API og en frontend bygget med OpenLayers for \textbf{kartbasert visualisering}. Tilbakemeldinger fra interessenter bekrefter at prototypen er gjennomførbar og effektiv som grunnlag for et fremtidig operativt system i skogbruksnæringen.
    } \\ \hline
  \end{tabularx}
\end{table}
